% Gerado por lcqui-spec-doc; não editar.
\subsection*{Emprestimo\_Reagente --- projeção M3 completa (34 colunas)}
Registro relacional completo do empréstimo (34 colunas), incluindo o snapshot imutável vencido\_na\_retirada. Não é payload de criação nem documento Firestore completo. A existência global das FKs, a unicidade de empréstimo ativo por frasco e as transições de status são verificadas fora do registro (loan\_state.als).
\begin{description}
\item[id] Tipo: int. Obrigatório: sim. Aceita nulo: não.
SQL: SERIAL.
PK; geração e unicidade globais fora do registro local.
\item[id\_frasco\_reagente] Tipo: int. Obrigatório: sim. Aceita nulo: não.
SQL: INTEGER.
FK para Frasco\_Reagente; existência e unicidade ativa validadas fora do registro.
\item[id\_usuario\_retirou] Tipo: int. Obrigatório: sim. Aceita nulo: não.
SQL: INTEGER.
Portador físico (Professor/Bolsista) que recebe e usa o reagente; FK verificada pelo servidor.
\item[id\_gestor\_retirada] Tipo: int. Obrigatório: sim. Aceita nulo: não.
SQL: INTEGER.
Operador do almoxarifado que registrou a retirada; no autoatendimento coincide com o retirante.
\item[status] Tipo: enum. Obrigatório: sim. Aceita nulo: não.
SQL: ENUM.
Valores: EM\_USO, ATRASADO, DEVOLVIDO, DEVOLVIDO\_COM\_ATRASO, DEVOLVIDO\_COM\_ANOMALIA, ENCERRADO\_EXTRAORDINARIO.
Estado controlado pelo servidor; EM\_USO/ATRASADO são custódia ativa.
\item[data\_retirada] Tipo: string. Obrigatório: sim. Aceita nulo: não.
SQL: TIMESTAMP.
\item[data\_devolucao\_prevista] Tipo: string. Obrigatório: sim. Aceita nulo: não.
SQL: DATE.
Padrão: \textasciicircum{}[0-9]\{4\}-[0-9]\{2\}-[0-9]\{2\}\$.
Data civil YYYY-MM-DD; vencimento legal às 23:59:59.999 em America/Sao\_Paulo.
\item[data\_devolucao\_efetuada] Tipo: string. Obrigatório: sim. Aceita nulo: sim.
SQL: TIMESTAMP.
\item[id\_local\_usado] Tipo: int. Obrigatório: sim. Aceita nulo: não.
SQL: INTEGER.
\item[id\_usuario\_devolveu] Tipo: int. Obrigatório: sim. Aceita nulo: sim.
SQL: INTEGER.
\item[id\_gestor\_devolucao] Tipo: int. Obrigatório: sim. Aceita nulo: sim.
SQL: INTEGER.
\item[peso\_saida] Tipo: number. Obrigatório: sim. Aceita nulo: não.
SQL: NUMERIC(10,3).
Mínimo: 0.
Máximo: 9999999.999.
Múltiplo de 0.001.
Peso bruto em g medido na retirada; referência imutável do consumo.
\item[peso\_retorno] Tipo: number. Obrigatório: sim. Aceita nulo: sim.
SQL: NUMERIC(10,3).
Mínimo: 0.
Máximo: 9999999.999.
Múltiplo de 0.001.
Peso bruto em g medido na devolução; preservado mesmo quando contradiz tara/saída.
\item[medida\_utilizada] Tipo: number. Obrigatório: sim. Aceita nulo: sim.
SQL: NUMERIC(10,3).
Mínimo: 0.
Máximo: 9999999.999.
Múltiplo de 0.001.
Consumo validado não negativo na unidade do empréstimo; null até a validação quantitativa. Zero é consumo válido.
\item[consumo\_validado] Tipo: bool. Obrigatório: sim. Aceita nulo: não.
SQL: BOOLEAN.
Falso na retirada e enquanto a pendência quantitativa permanece aberta; server-owned.
\item[anomalia\_metrologica] Tipo: string. Obrigatório: sim. Aceita nulo: sim.
SQL: TEXT.
Classificação server-owned da anomalia quantitativa; null quando não há anomalia.
\item[peso\_retorno\_efetivo] Tipo: number. Obrigatório: sim. Aceita nulo: sim.
SQL: NUMERIC(10,3).
Mínimo: 0.
Máximo: 9999999.999.
Múltiplo de 0.001.
Peso de retorno efetivo em g; null enquanto o consumo não for validado.
\item[id\_resolucao\_metrologica] Tipo: int. Obrigatório: sim. Aceita nulo: sim.
SQL: INTEGER.
Referência ao evento AJUSTE que encerrou a pendência; resolver não apaga o status anômalo.
\item[unidade\_medida\_utilizada] Tipo: enum. Obrigatório: sim. Aceita nulo: não.
SQL: ENUM.
Valores: ml, g.
Qualifica exclusivamente medida\_utilizada: ml para líquido, g para sólido.
\item[densidade\_aplicada] Tipo: number. Obrigatório: sim. Aceita nulo: sim.
SQL: NUMERIC(10,5).
Maior que 0.
Máximo: 99999.99999.
Múltiplo de 0.00001.
Snapshot histórico da densidade na retirada; obrigatório e positivo para líquido, null para sólido.
\item[peso\_perda\_evaporacao] Tipo: number. Obrigatório: sim. Aceita nulo: não.
SQL: NUMERIC(10,3).
Mínimo: 0.
Máximo: 9999999.999.
Múltiplo de 0.001.
Perda em g registrada separadamente do consumo; não inferir arbitrariamente.
\item[uso\_vencido\_aceito] Tipo: bool. Obrigatório: sim. Aceita nulo: não.
SQL: BOOLEAN.
Ciência explícita do empréstimo específico; não substitui autorização do gestor.
\item[vencido\_na\_retirada] Tipo: bool. Obrigatório: sim. Aceita nulo: não.
SQL: BOOLEAN.
Snapshot server-owned e imutável do vencimento conhecido no instante em que a retirada cria o empréstimo, após eventual primeira abertura. Não é autorização nem é recalculado na devolução.
\item[finalidade\_uso] Tipo: enum. Obrigatório: sim. Aceita nulo: não.
SQL: ENUM.
Valores: AULA\_PRATICA, DEMONSTRACAO, PESQUISA\_TCC\_POS, ESTUDO\_DEGRADACAO\_RESIDUOS.
Finalidade declarada na retirada; pesquisa excepcional exige TCR de Q04.
\item[auto\_atendimento] Tipo: bool. Obrigatório: sim. Aceita nulo: não.
SQL: BOOLEAN.
Calculado pelo backend conforme Q14; justificativa e notificação à chefia quando true.
\item[justificativa\_metodologica] Tipo: string. Obrigatório: sim. Aceita nulo: sim.
SQL: TEXT.
Condicional ao uso excepcional de pesquisa Q04 (mínimo verificado após trim).
\item[tcr\_versao] Tipo: string. Obrigatório: sim. Aceita nulo: sim.
SQL: VARCHAR(100).
Máximo: 100 caracteres.
\item[tcr\_aceito\_em] Tipo: string. Obrigatório: sim. Aceita nulo: sim.
SQL: TIMESTAMP.
\item[tcr\_aceito\_por] Tipo: int. Obrigatório: sim. Aceita nulo: sim.
SQL: INTEGER.
\item[tcr\_auditoria\_id] Tipo: int. Obrigatório: sim. Aceita nulo: sim.
SQL: INTEGER.
\item[tipo\_encerramento\_excepcional] Tipo: enum. Obrigatório: sim. Aceita nulo: sim.
SQL: ENUM.
Valores: QUEBRA\_ACIDENTAL, EXTRAVIO\_SINISTRO.
Tipo do sinistro; null fora do encerramento extraordinário.
\item[motivo\_encerramento\_excepcional] Tipo: string. Obrigatório: sim. Aceita nulo: sim.
SQL: TEXT.
\item[id\_gestor\_encerramento] Tipo: int. Obrigatório: sim. Aceita nulo: sim.
SQL: INTEGER.
\item[massa\_perda\_estimada\_g] Tipo: number. Obrigatório: sim. Aceita nulo: sim.
SQL: NUMERIC(10,3).
Mínimo: 0.
Máximo: 9999999.999.
Múltiplo de 0.001.
\end{description}
